\subsection{Accuracy on multi-turn conversations}\label{sec:accuracy}
All four modes answer 80--83 of the 84 questions (Table~\ref{tab:accuracy}). No mode is clearly worse than exact routing.

\begin{table}[t]
\centering\small
\caption{Conversation suite: 9 Arabic conversations, 84 questions, 107 automatic checks. Speed is the mean writing speed over all answers.}
\label{tab:accuracy}
\begin{tabular}{lccc}
\toprule
Mode & Questions passed & Checks passed & Speed (tokens/s) \\
\midrule
off (exact) & 80 / 84 & 103 / 107 & 4.77 \\
light & 83 / 84 & 106 / 107 & 5.39 \\
medium & 80 / 84 & 103 / 107 & 5.73 \\
strong & 83 / 84 & 106 / 107 & 6.02 \\
\bottomrule
\end{tabular}
\end{table}

The exact and light numbers combine the first full run with a rerun of the hospital-report conversation. That rerun came after a fix that bars the model from opening a thinking block when thinking is off. The fix changes no other turn, because no other turn opened a thinking block.

The failures are ordinary model errors, not routing artifacts:
\begin{itemize}
  \item \textbf{Every mode} wrote 6 words when asked for exactly 5.
  \item \textbf{Exact routing} forgot a 600,000-dinar expense in the ledger, confused grandfather and grandmother in the family conversation, and counted 9 letters in a 10-letter word.
  \item \textbf{Medium routing} repeated the family error, ordered two historical events wrongly despite giving the correct years, and ran out of tokens in the rota without a final answer line.
\end{itemize}

The ranking of the modes (light = strong $>$ off = medium) does not follow routing strength, which points to noise. A three-question difference out of 84 cannot separate the modes. Section~\ref{sec:fidelity} measures what this test cannot.

\subsection{Fidelity of routing}\label{sec:fidelity}
The routed modes keep the exact model's top choice on 98.9--99.2\% of tokens (Table~\ref{tab:fidelity}, first run). Every cost of routing grows with its strength, as the speed gain does.

\begin{table}[t]
\centering\small
\caption{Fidelity test: 10,840 answer tokens of the exact model, read back under each routing mode. KL is measured over the reference's 32 likeliest tokens plus one bucket for the rest.}
\label{tab:fidelity}
\setlength{\tabcolsep}{4pt}
\begin{tabular}{lccccccc}
\toprule
Mode & \makecell{Same top\\token} & \makecell{Mean KL\\(nats/token)} & \makecell{Tokens with\\KL $>$ 0.1} & \makecell{Perplexity\\change} & \makecell{Speed\\(tokens/s)} & \makecell{Disk read\\per token} & \makecell{Cache\\hits} \\
\midrule
off (reference) & 100\% & 0 & 0\% & --- & 4.61 & 0.86 GB & 82.7\% \\
light & 99.20\% & 0.0036 & 0.77\% & +0.24\% & 5.25 & 0.71 GB & 87.6\% \\
medium & 98.94\% & 0.0048 & 1.13\% & +0.24\% & 5.59 & 0.64 GB & 89.6\% \\
strong & 98.97\% & 0.0052 & 1.22\% & +0.29\% & 5.80 & 0.60 GB & 90.8\% \\
\bottomrule
\end{tabular}
\end{table}

Each mode read the same 10,840 answer tokens in 84 turns of 9 conversations. Speed here is the speed of reading answers one token at a time, a close proxy for writing them. Strong routing cuts disk traffic per token by 30\% and wait time on the disk from 802\,s to 355\,s over the run.

\paragraph{Repeatability.} We ran the test a second time for exact, light and strong routing (Table~\ref{tab:repeat}). Exact routing reproduced the first run bit for bit on all 10,840 tokens: the same top choice, summed log-likelihood and agreement with the stored text in every turn. Routed runs depend on what the disk has delivered by each step, so they vary a little, but every measure keeps the modes in the same order. The perplexity change is the least stable measure: light routing moved from +0.24\% to +0.15\% between runs. Agreement and KL move by far less than the gap between the modes.

\begin{table}[t]
\centering\small
\caption{Two runs of the fidelity test for the routed modes. Exact routing was identical in both.}
\label{tab:repeat}
\begin{tabular}{lccccc}
\toprule
Mode & \makecell{Same top\\token} & \makecell{Mean KL\\(nats/token)} & \makecell{Tokens with\\KL $>$ 0.1} & \makecell{Perplexity\\change} & \makecell{Speed\\(tokens/s)} \\
\midrule
light, run 1 & 99.20\% & 0.0036 & 0.77\% & +0.24\% & 5.25 \\
light, run 2 & 99.06\% & 0.0037 & 0.83\% & +0.15\% & 5.14 \\
strong, run 1 & 98.97\% & 0.0052 & 1.22\% & +0.29\% & 5.80 \\
strong, run 2 & 98.93\% & 0.0053 & 1.34\% & +0.32\% & 5.67 \\
\bottomrule
\end{tabular}
\end{table}

Agreement varies little across conversations. JSON is least affected: 99.9--100\% for every mode. The Baghdad history conversation is most affected under medium and strong routing (97.9\% and 98.1\%). The worst single turns still agree on 90--96\% of tokens across both runs.

Under exact routing, the model's top token matches the stored answer text on 99.2\% of positions, not 100\%. Re-tokenizing generated text does not always reproduce the tokens the model wrote. This affects all modes equally, because all compare against the same exact distributions.

\section{Discussion and limitations}
The expert store gives most of the gain; routing adds a smaller, tunable part. Over 512 written tokens, exact routing already turns llama.cpp's 3.7 tokens/s into 5.5, and strong routing adds another 19\%; prompts are read 6--7$\times$ faster. A user who wants the exact model loses nothing by using Rafid.

\paragraph{Accuracy tests cannot rank routing modes at this scale.} On 84 questions, the four modes land within three answers of each other, in an order unrelated to routing strength. The fidelity test does rank them, consistently on every measure and in both runs. We suggest that work on approximate routing report token-level fidelity alongside task scores.

\paragraph{One token in a hundred.} Strong routing changes the exact model's top choice on about 1 token in 100, and clearly shifts its distribution on 1.2\%. In free generation, the first changed token sends the text on a different path. So outputs differ more than the per-token number suggests. Of the 84 answers, 37 under light routing were word-for-word identical to the exact ones, 27 under medium and 23 under strong. Whether those different paths are worse is what accuracy tests measure, and there we saw no loss.

\paragraph{Measurements still missing.}
\begin{itemize}
  \item \textbf{Standard benchmarks.} Our questions are Arabic and conversational. Perplexity on a standard corpus and a GSM8K or MMLU subset with free-form reasoning would allow comparison with prior work. Answers must be written, not chosen from the first token, because routing only acts while writing.
  \item \textbf{Other models and machines.} We measured one model on one laptop. Rafid also runs Qwen3.5-122B-A10B and Qwen3.6-35B-A3B, which have not been evaluated here.
  \item \textbf{A direct comparison with cache-prior routing.} The bias rule of \citet{skliar2024} could be implemented in the same store and compared at matched hit rates.
\end{itemize}

\paragraph{Threats to validity.} The reference answers for 8 of the 9 conversations come from a run before the thinking-tag fix, and the report conversation from a run after it. Both used exact routing and the same model, and no answer in those 8 conversations opened a thinking block. The checks are automatic: ``any mention of the number'' checks can accept a right number reached for the wrong reason. Each turn's example answer allows a manual review.

\section{Conclusion}
A 397B-parameter MoE model runs on a 31 GB Windows laptop. With exact routing it writes 5.5 tokens/s over 512 tokens, 1.5$\times$ llama.cpp's memory-mapped baseline, and reads prompts 6--7$\times$ faster. It writes 6.5 tokens/s, 1.75$\times$, when the model may swap a low-ranked missing expert for a cached one. That swap changes about one token choice in a hundred and no measured answer quality. Token-level fidelity, measured inside real conversations, is what makes the cost of approximate routing visible, where answer counts are blind to it.

\section*{Code and data availability}
Rafid is released under the MIT licence: source, a Windows binary build, the Arabic conversation suite, an English template suite, and the raw result files of every table (\texttt{docs/results}, with \texttt{docs/REPRODUCE.md}). \texttt{tools/summarize\_results.py} rebuilds every table from those files. Evaluators can import their own suites; each run records the SHA-256 of the suite file it used. Repository: \url{https://github.com/zaidlatif2010/rafid}.

\begin{thebibliography}{6}
\bibitem[Alizadeh et al.(2023)]{alizadeh2023}
K.~Alizadeh et al.
\newblock LLM in a flash: Efficient large language model inference with limited memory.
\newblock arXiv:2312.11514, 2023.

\bibitem[Eliseev and Mazur(2023)]{eliseev2023}
A.~Eliseev and D.~Mazur.
\newblock Fast inference of mixture-of-experts language models with offloading.
\newblock arXiv:2312.17238, 2023.

\bibitem[Tang et al.(2024)]{tang2024}
P.~Tang et al.
\newblock HOBBIT: A mixed precision expert offloading system for fast MoE inference.
\newblock arXiv:2411.01433, 2024.

\bibitem[Skliar et al.(2024)]{skliar2024}
A.~Skliar et al.
\newblock Mixture of cache-conditional experts for efficient mobile device inference.
\newblock arXiv:2412.00099, 2024.

\bibitem[Woods(2026)]{woods2026}
D.~Woods.
\newblock Flash-MoE.
\newblock As reported in Simon Willison's weblog, \url{https://simonwillison.net/b/9337}, March 2026; code at \url{https://github.com/fiveangle/flash-moe} (frozen fork).

\bibitem[Gerganov et al.(2023)]{llamacpp}
G.~Gerganov and contributors.
\newblock llama.cpp.
\newblock \url{https://github.com/ggml-org/llama.cpp}, builds b12167 and b12169.
\end{thebibliography}
